\documentclass[11pt]{article}
\usepackage[margin=0.8in]{geometry}
\usepackage{amsmath,amssymb,booktabs,array,longtable}
\usepackage{hyperref}
\usepackage{enumitem}
\usepackage{graphicx}
\usepackage{float}
\setlist{nosep}

\begin{document}

\begin{center}
{\Large\textbf{Shiv Nadar University Chennai}}\\
\textbf{CS3807 -- Deep Learning Laboratory}\\
\textbf{Experiment 2}\\
{\large\textbf{Implementation of a Multi-Layer Perceptron (MLP) for Multi-Class Image Classification}}
\end{center}

\renewcommand{\arraystretch}{1.25}

\begin{tabular}{|p{3.8cm}|p{8cm}|p{2cm}|}
\hline
Degree \& Branch & B.Tech Artificial Intelligence \& Data Science & Semester V\\ \hline
Subject Code \& Name & CS3807 -- Deep Learning Laboratory & AY: 2026--27\\ \hline
\end{tabular}

\section*{1. Objective}
Implement an MLP using TensorFlow/Keras on the Fashion-MNIST dataset. Learn image preprocessing, flattening, model construction, training, evaluation, and automated hyperparameter optimization.

\section*{2. Learning Outcomes}
Students will be able to:
\begin{enumerate}
\item Explain the architecture of an MLP.
\item Preprocess image datasets.
\item Build and train an MLP.
\item Evaluate classification performance.
\item Perform automated hyperparameter optimization.
\item Recommend the best model configuration.
\end{enumerate}

\section*{3. Background Theory}
An MLP consists of an input layer, hidden layers and an output layer.

\[
z^{(l)}=W^{(l)}a^{(l-1)}+b^{(l)},\qquad
a^{(l)}=f(z^{(l)})
\]

Output layer:
\[
\hat y=\mathrm{Softmax}(z)
\]

Loss:
\[
L=-\sum_i y_i\log(\hat y_i)
\]

Recommended architecture:

\begin{center}
784 $\rightarrow$ Dense(128, ReLU) $\rightarrow$ Dense(64, ReLU) $\rightarrow$ Dense(10, Softmax)
\end{center}

\section*{4. Dataset}
Dataset: Fashion-MNIST

Training Images: 60,000 \\
Testing Images: 10,000 \\
Classes: 10 \\
Image Size: $28\times28$

\section*{5. Image Preprocessing}

Fashion-MNIST images are stored as $28\times28$ matrices. Dense layers require a one-dimensional feature vector.

\[
28\times28=784
\]

Example:

\[
\begin{bmatrix}
10&20\\30&40
\end{bmatrix}
\rightarrow
[10,20,30,40]
\]

Perform the following preprocessing:

\begin{enumerate}
\item Load Fashion-MNIST.
\item Display sample images.
\item Flatten images.
\item Normalize pixels to [0,1].
\item Convert labels into one-hot vectors.
\end{enumerate}

\section*{6. Experimental Procedure}

\textbf{Task 1: Dataset Exploration}
\begin{itemize}
\item Load Fashion-MNIST.
\item Print dataset dimensions.
\item Display ten sample images.
\item Plot class distribution.
\end{itemize}

\textbf{Output obtained:}
\begin{itemize}
\item Raw Training Data Shape: (60000, 28, 28)
\item Raw Testing Data Shape: (10000, 28, 28)
\end{itemize}

\bigskip

\textbf{Task 2: Data Preprocessing}

\begin{itemize}
\item Flatten images.
\item Normalize pixel values.
\item Print tensor shapes before and after preprocessing.
\end{itemize}

\textbf{Output obtained:}
\begin{itemize}
\item Preprocessed Training Data Shape: (60000, 784)
\item Preprocessed Testing Data Shape: (10000, 784)
\end{itemize}

\bigskip

\textbf{Task 3: Model Construction}

Construct an MLP with

\begin{itemize}
\item Input layer (784)
\item Dense(128, ReLU)
\item Dense(64, ReLU)
\item Dense(10, Softmax)
\end{itemize}

\texttt{model.summary()} output: Total params 109,386 (427.29 KB) -- all trainable.

\bigskip

\textbf{Task 4: Model Training}

Compile using

\begin{itemize}
\item Optimizer: Adam
\item Loss: Categorical Cross Entropy
\item Metric: Accuracy
\end{itemize}

Trained for 20 epochs, batch size 32. Final epoch (20/20): training accuracy 0.9392, training loss 0.1600, validation accuracy 0.8826, validation loss 0.4132.

\bigskip

\textbf{Task 5: Model Evaluation}

Baseline test accuracy: \textbf{0.8805}. Full precision/recall/F1 per class obtained via \texttt{classification\_report}; macro avg precision/recall/F1 all $\approx 0.88$. See Performance Comparison table in Section 9 for consolidated figures.

\section*{7. Hyperparameter Optimization}

Instead of manually selecting hyperparameters, students shall perform automated hyperparameter optimization using either \textbf{GridSearchCV} or \textbf{RandomizedSearchCV} (recommended) together with the SciKeras wrapper.

Optimize the following search space.

\begin{center}
\begin{tabular}{|p{5cm}|p{7cm}|}
\hline
Hyperparameter & Candidate Values\\
\hline
Hidden Layers & 1, 2, 3\\
Hidden Neurons & 32, 64, 128, 256\\
Learning Rate & 0.1, 0.01, 0.001\\
Batch Size & 16, 32, 64, 128\\
Epochs & 10, 20, 30\\
Optimizer & SGD, Adam, RMSProp\\
Activation Function & ReLU, Tanh, Sigmoid\\
Dropout Rate (Optional) & 0.0, 0.2, 0.5\\
\hline
\end{tabular}
\end{center}

\subsection*{Tasks}

\begin{enumerate}
\item Build a baseline MLP model.
\item Define the hyperparameter search space.
\item Perform Grid Search or Randomized Search using 5-fold cross-validation.
\item Record the best hyperparameter combination.
\item Retrain the model using the optimized hyperparameters.
\item Evaluate the optimized model on the testing dataset.
\item Compare the optimized model with the baseline model.
\end{enumerate}

\textit{Note: the run used RandomizedSearchCV with \textbf{3-fold} cross-validation and 5 candidate configurations (a narrowed search, scaled down from the manual's 5-fold spec for execution-time reasons). If your grading criteria require literal 5-fold CV, re-run the search with \texttt{cv=5} before submitting.}

\section*{8. Mandatory Plots}

Generate the following plots. Write a 2--3 line inference underneath each.

\textbf{1. Sample Images}
\begin{figure}[H]
\centering
\includegraphics[width=0.85\textwidth]{figs/plot1.png}
\end{figure}
\textit{The 10 sample images span the various types of garments observed in the dataset for analysis.}

\bigskip
\textbf{2. Class Distribution}
\begin{figure}[H]
\centering
\includegraphics[width=0.75\textwidth]{figs/plot2.png}
\end{figure}
\textit{Each of the 10 classes gas 6000 training images, for sake of balance.}

\bigskip
\textbf{3 \& 4. Training Accuracy vs Epoch / Validation Accuracy vs Epoch}
\begin{figure}[H]
\centering
\includegraphics[width=0.7\textwidth]{figs/plot3_and_4.png}
\end{figure}
\textit{Accuracy climbs to 0.94 by the 20th iteration, while accuracy stagnates at 0.88 at epoch 5, after which it oscillates.}

\bigskip
\textbf{5 \& 6. Training Loss vs Epoch / Validation Loss vs Epoch}
\begin{figure}[H]
\centering
\includegraphics[width=0.7\textwidth]{figs/plot5_and_6.png}
\end{figure}
\textit{Baseline training loss falls to 0.16, validation loss reaches 0.34 at around epoch 9. It then rises to 0.42 which shows overfitting. The optimised model's validation loss still decreases.}

\bigskip
\textbf{7. Confusion Matrix}
\begin{figure}[H]
\centering
\includegraphics[width=0.85\textwidth]{figs/plot7.png}
\end{figure}
\textit{Shirt is a hard class to classify with only 617 correct ones. Pullover and coat are somewhere in the middle. Trousers, bag, etc are classified almost perfectly.}

\bigskip
\textbf{8. Hyperparameter Search Results}
\begin{figure}[H]
\centering
\includegraphics[width=0.85\textwidth]{figs/plot8.png}
\end{figure}
\textit{Mean CV accuracy across the 5 configurations ranges from 0.7 to 0.85. 1 layer. 256 neurons, tanh, Adam, learning rate 0.001 was found to be the best.}

\bigskip
\textbf{9. Best Model Accuracy Comparison}
\begin{figure}[H]
\centering
\includegraphics[width=0.6\textwidth]{figs/plot9.png}
\end{figure}
\textit{Optimised model scores 0.8754 compared to baseline model's 0.8805 i.e. a faster model was found but it was slightly worse than the baseline model. The time saved is significant though.}

\section*{9. Results}

\textbf{Best Hyperparameters}

\begin{tabular}{|p{6cm}|p{6cm}|}
\hline
Hidden Layers & 1\\\hline
Hidden Neurons & 256\\\hline
Learning Rate & 0.001\\\hline
Batch Size & 64\\\hline
Optimizer & Adam\\\hline
Activation Function & Tanh\\\hline
Epochs & 10\\\hline
Dropout & 0.0\\\hline
Cross-validation Accuracy & 0.8473\\\hline
Testing Accuracy & 0.8754\\\hline
\end{tabular}

\vspace{3mm}

\textbf{Performance Comparison}

\begin{tabular}{|p{4cm}|c|c|}
\hline
Metric & Baseline & Optimized\\
\hline
Accuracy & 0.8805 & 0.8754\\
Precision & 0.8831 & 0.8770\\
Recall & 0.8805 & 0.8754\\
F1-score & 0.8811 & 0.8745\\
Training Time & 162.78s & 59.28s\\
\hline
\end{tabular}

\section*{10. Discussion}
 
\begin{enumerate}
\item \textbf{Which optimization method was used?} \\
\textit{RandomizedSearchCV with 3 fold cross validation over 5 sampled configurations}
\item \textbf{Which hyperparameters were selected?} \\
\textit{1 hidden layer, 256 neurons, learning rate 0.001, batch size 64, Adam optimizer, Tanh activation, 10 epochs.}
\item \textbf{Did optimization improve performance?} \\
\textit{The optimized model scored 0.8754 vs the baseline's 0.8805 test accuracy, a slight decrease. However, what improved was training efficiency: 59.28s vs 162.78s, significantly faster}
\item \textbf{Which hyperparameter had the greatest impact?} \\
\textit{Using only 1 hidden layer (vs the baseline's 2) and still landing within half a point of baseline accuracy, hidden layer depth appears to trade off against training time more than it costs accuracy for this dataset.}
\item \textbf{Compare Grid Search and Randomized Search.} \\
\textit{Grid Search exhaustively evaluates every combination in the search space, guaranteeing the best result within the grid. Randomized Search samples a fixed number of combinations at random. It is less exhaustive but takes a lot less time.}
\item \textbf{Recommend the best MLP configuration.} \\
\textit{1 hidden layer of 256 neurons with Tanh activation, Adam optimizer at learning rate 0.001, batch size 64, trained for 10 epochs with no dropout.}
\end{enumerate}



\section*{11. References}
\begin{enumerate}
\item Goodfellow et al., \emph{Deep Learning}.
\item Bishop, \emph{Pattern Recognition and Machine Learning}.
\item Haykin, \emph{Neural Networks and Learning Machines}.
\item Fashion-MNIST Dataset.
\item TensorFlow/Keras Documentation.
\end{enumerate}

\end{document}